\section{Introduction}

Understanding the evolutionary trajectory of the human brain is a central goal of comparative neuroscience, because the adaptations that distinguish \textit{Homo sapiens} from other primates underpin the cognitive abilities that make our species unique \cite{rilling2014comparative, sherwood2013whats}. Rhesus macaques, our closest cercopithecoid relatives, share extensive neuroanatomical homology with humans and have therefore served as the reference species against which human-specific brain architecture is benchmarked. By leveraging magnetic resonance imaging (MRI) in both species, investigators have documented reorganizations of the cortical sheet \cite{hill2010similar, smaers2019brain}, specializations of language-related white-matter tracts \cite{rilling2008evolution, eichert2019special, balezeau2020primate}, and differences in macroscale structural connectivity \cite{goulas2014comparative}. Yet a foundational question remains open: does the human brain only differ from the macaque brain in where anatomy is arranged, or also in \textit{when} that anatomy matures along development?

Two complementary bodies of work motivate this question. First, \textit{spatial} comparative neuroimaging has exposed a human-specific expansion of the arcuate fasciculus (AF), a dorsal language pathway that projects more prominently into the middle and inferior temporal lobes of humans than of chimpanzees or macaques \cite{rilling2008evolution, rilling2012continuity, catani2008arcuate}. Cortical surface-based analyses similarly identify regions of the temporoparietal junction, posterior superior temporal sulcus, and prefrontal cortex that are disproportionately expanded or reconfigured in humans \cite{hill2010similar, patel2019temporoparietal, smaers2011primate}. Second, \textit{temporal} studies of brain development have shown that human neocortical myelination unfolds over a protracted adolescence that extends years beyond the chimpanzee trajectory \cite{miller2012prolonged, perrin2009sex, dubois2014early}, while rhesus macaques show comparatively compressed pubertal maturation. These parallel literatures converge on the suggestion that anatomical and developmental differences are entangled, but they have rarely been analyzed in the same quantitative framework.

Machine-learning models that predict chronological age from structural MRI provide a promising way to collapse anatomy onto a single developmental axis. Brain-age prediction has matured rapidly in humans, with biomarkers such as the brain-age gap (the difference between predicted and chronological age) now linked to psychiatric illness \cite{cole2017predicting, franke2010estimating, hajek2019brain, mcwhinney2021obesity}. However, brain-age models have almost exclusively been developed and evaluated \textit{within} a single species; they have not been repurposed to ask whether the anatomical features that encode chronological age in one primate species also encode age in another, nor whether cross-species prediction errors can themselves serve as an index of evolutionary divergence.

Here we close this gap. We combined structural and diffusion MRI from 370 humans aged 8--14 years (Human Connectome Project--Development \cite{somerville2018lifespan, harms2018extending}) and 181 rhesus macaques aged 2--4 years to build brain structure based cross-species age prediction models. Gray-matter volume (GMV) and four diffusion tensor imaging (DTI) indices -- fractional anisotropy, mean diffusivity, axial diffusivity, and radial diffusivity -- were extracted over 92 regional-map parcels \cite{bezgin2012hierarchical} and 42 XTRACT-defined white-matter tracts \cite{warrington2020xtract}, yielding 260 candidate features per subject. After ComBat harmonization to compensate for scanner heterogeneity \cite{fortin2017harmonization_dti, fortin2017harmonization_thickness, johnson2007adjusting} and INIA19-guided spatial normalization for macaques \cite{rohlfing2012inia19}, we trained within-species age predictors with 10-fold cross-validation and then applied them across species. From the cross-species predictions we defined a \textit{brain cross-species age gap} (BCAP) as a percentile-rank difference, and correlated BCAP with behavioral phenotypes, gray-matter regions, and white-matter tracts to localize evolutionary divergence along the chronological axis.

\paragraph{Contributions.}
\begin{itemize}
  \item We introduce the first brain structure based cross-species age prediction framework that jointly characterizes human and macaque brain anatomy on a shared chronological axis.
  \item We demonstrate a stable \textit{asymmetry} in cross-species prediction: the trained macaque model predicts human age with higher accuracy (R = 0.4823, MAE = 8.3610 years) than the trained human model predicts macaque age (R = 0.2898, MAE = 7.6157 years).
  \item We show that prediction error scales with actual age, in opposite directions for the two species, consistent with a growing evolutionary divergence across human adolescence.
  \item We propose BCAP, a percentile-rank index that quantifies cross-species developmental divergence per subject, and link it to both language (picture vocabulary test: R = 0.1588) and perceptual (visual sensitivity test: R = -0.2051) behavior.
  \item We map evolution-associated structure, localizing positive BCAP weights to the arcuate fasciculus and right prefrontal cortex and negative weights to putamen, amygdala, and subcortical nuclei.
\end{itemize}

\section{Related Work}

\paragraph{Spatial comparative neuroanatomy of humans and macaques.}
Voxel-based and surface-based analyses have established that the human cortical sheet is disproportionately expanded over lateral prefrontal, temporoparietal, and parietal association regions relative to macaques \cite{hill2010similar, smaers2011primate, smaers2019brain}. The arcuate fasciculus, which connects inferior frontal and posterior temporal language areas, shows a marked leftward lateralization and expanded temporal projection in humans that is absent or reduced in macaques and attenuated in chimpanzees \cite{rilling2008evolution, rilling2012continuity, eichert2019special}. Auditory-cortex projections of the AF appear to have evolved from a nonhuman primate prototype, supporting a continuity-with-divergence account of language pathway evolution \cite{balezeau2020primate, catani2008arcuate, rocchi2021common}. These spatial findings anchor our selection of tracts (AF, IFOF, SLF, CT, UF) and regions (prefrontal cortex, insula, basal ganglia) as candidate features, but they do not themselves quantify how quickly these structures mature in each species.

\paragraph{Developmental neuroimaging in humans and nonhuman primates.}
Longitudinal and cross-sectional studies in humans show slow, prolonged white-matter maturation extending into the third decade, with sex- and tract-specific trajectories \cite{dubois2014early, perrin2009sex, miller2012prolonged}. Macaque brain development follows a comparable ordering but on a compressed timescale, with sensorimotor cortex maturing earlier than association cortex. Our use of HCP-Development \cite{somerville2018lifespan, harms2018extending} for 8--14 year-old humans and a 2--4 year-old macaque cohort is motivated by approximate age-range matching of these developmental windows. By building an age-prediction model on top of these cohorts, we shift the analytic focus from ``where'' the species differ to ``how far apart'' their developmental trajectories are at each chronological point.

\paragraph{Brain-age prediction and the brain-age gap.}
Machine-learning brain-age models estimate chronological age from structural MRI features, and the residual -- the brain-age gap -- has emerged as a biomarker of atypical maturation in schizophrenia, bipolar disorder, and obesity-related aging \cite{cole2017predicting, franke2010estimating, hajek2019brain, mcwhinney2021obesity}. To date, however, these models have been trained and evaluated entirely within one species. We extend this paradigm by deliberately applying within-species predictors \textit{across} species and interpreting the resulting errors as a signature of evolutionary divergence. The percentile-rank BCAP index we propose avoids the scale mismatch between human and macaque ages that would afflict raw residuals and makes the cross-species discrepancy interpretable at the individual level.

\paragraph{Cross-species atlases and scanner harmonization.}
Reliable cross-species feature extraction requires homologous parcellations and modality-matched templates. The XTRACT protocol \cite{warrington2020xtract} provides 42 white-matter tracts defined consistently across human and macaque diffusion data, and the INIA19 template \cite{rohlfing2012inia19} together with the regional-map gray-matter parcellation \cite{bezgin2012hierarchical} enables gray-matter volume extraction in a homologous frame. Because macaque MRI was acquired on two GE 3.0T scanners with different sequences, we also applied ComBat harmonization \cite{johnson2007adjusting, fortin2017harmonization_dti, fortin2017harmonization_thickness} to mitigate scanner-induced variance. Together these choices ensure that the cross-species prediction asymmetry we report is not confounded by atlas or site effects.
